\section{Future Work}
\label{sec:future_work}

The system was built and operated under tournament time pressure, and
several extensions were deliberately deferred. We collect here the ones
we consider most valuable: richer data sources, upgrades to the score
model and team ratings, principled Bayesian updating, optimizer
extensions, and the prediction-market combiner that we regard as the
main methodological follow-up.

\subsection{Data Sources}
\label{sec:future-data}

The single largest gap during live operation was confirmed starting
elevens. A scraper that monitors manager press conferences, predicted
lineups from major outlets, official team sheets released an hour
before kickoff, and national-team beat reporters, and that emits a
per-player starting probability before each deadline, would have
closed the failure mode that cost us most: rotation by teams that had
already clinched their group.

A source audit conducted during the tournament identified four
additions worth integrating. The \texttt{soccerdata} Python library
wraps eight providers behind one interface and would add shot-level
expected goals, lineups, and historical bookmaker odds useful for
backtest calibration, though its shot-level coverage is limited to
club leagues. The \texttt{worldfootballR} R package is the only
verified route to FBref international coverage, including shot events
and advanced per-match statistics for past World Cups; the hard caveat
is that the package was archived in September 2025 and is unmaintained,
so it is suitable for historical feature building but not live
ingestion. StatsBomb Open Data provides free event data with per-shot
expected goals for the 2018 and 2022 World Cups and recent continental
tournaments, a natural independent validation set, although its user
agreement requires attribution and precludes redistributing the raw
data. Finally, the Wyscout event dataset of
\citet{pappalardo2019public} covers the 2018 World Cup under a
genuinely open license; deriving historical fantasy points from its
events would let us backtest the expected-points pipeline on a real
World Cup rather than on club proxies.

\subsection{Score-Model and Rating Upgrades}
\label{sec:future-methods}

Four method upgrades are cheap to implement and directly testable on
the existing walk-forward backtester.

\paragraph{Nested Poisson regression.} \citet{gilch2018elo} condition
the weaker team's goal rate on the stronger team's realized goals and
report that this nested variant backtested best among four Poisson
formulations on the 2010 and 2014 World Cups. It is a few-line change
to our Poisson backend and a natural next step beyond the independent
double-Poisson lineage \citep{maher1982modelling,karlis2003analysis}.

\paragraph{Dynamic Elo inside simulations.} The same study finds that
tournament simulations which freeze Elo ratings for the whole
tournament shift stage probabilities by up to five percentage points
and systematically favor stronger teams. Our Monte Carlo backend is
per-match with a fixed rating difference, so this applies once we
chain multi-round simulations; updating ratings \citep{elo1978rating}
between simulated rounds would recover underdog runs.

\paragraph{Dixon-Coles dependence.} Our simulator draws team and
opponent goals as independent Poisson variates. This independence
assumption is empirically misspecified: across 9{,}130 top-league
matches the home/away goal correlation is $-0.085$
\citep{petretta2025dependence}. The \citet{dixon1997modelling}
correction is the standard remedy, with a documented structural limit:
the adjustment only redistributes probability among the four
lowest-score cells, so Under/Over 2.5 probabilities are mathematically
identical to the independence model \citep{petretta2025dependence}.
Whether the correction moves fantasy expected points, as opposed to
match-outcome probabilities, is an open question that a walk-forward
A/B test answers cheaply; either outcome is informative, as a
justified upgrade or a justified deliberate simplification.

\paragraph{Pi-ratings.} The dynamic ratings of
\citet{constantinou2013pi}, updated from the discrepancy between
predicted and observed goal difference, outperformed football Elo
variants \citep{hvattum2010using} on league data. Their separate
home/away components matter less at the mostly neutral-venue World
Cup, so we scope this as a single ablation against our international
Elo feature.

Beyond these, a tournament-specific gradient-boosted model becomes
feasible now that the 2026 matches are recorded alongside prior World
Cups and continental tournaments. Per-tournament samples are small
(64 matches per World Cup), so heavy regularization is required to
avoid the overfitting we document earlier; the cross-domain transfer
problem, training on club data and predicting internationals, is
precisely what club-trained fantasy models such as OpenFPL
\citep{groos2025openfpl} do not face, and reducing that gap is the
point of the retraining.

\subsection{Bayesian Updating}
\label{sec:future-bayesian}

Our production blend of model prediction, realized points per match,
and recent form uses hand-tuned convex weights. The principled
replacement is a Bayesian update in which the model prediction is the
prior and realized match-by-match points are the likelihood, with the
data variance shrinking as matches accumulate. The fixed blend is an
acceptable approximation over a seven-match tournament; a full league
season would benefit from the proper treatment. The fuller version is
a hierarchical model with partial pooling: per-player skill, per-team
attack and defense strength, and per-manager rotation propensity, plus
an explicit minutes model in place of our fixed appearance assumption.
Such a model quantifies uncertainty as well as point estimates, which
the squad optimizer could consume as a chance-constrained variant of
the current deterministic program.

\subsection{Optimizer and Game-Mechanic Extensions}
\label{sec:future-optimizer}

The squad problem is at heart a salary-cap knapsack integer program
\citep{winston2022mathletics}; three refinements to our formulation
remain. First, captain selection should model the field explicitly:
the rank value of a captain choice depends on ownership and on the
distribution of the field's likely captains, and our current Monte
Carlo treats the field's captain as a single point estimate rather
than a distribution over candidates. Second, the group stage has a
multi-round horizon while knockouts are single games; a
stochastic-programming generalization would weight horizon-future
points by the probability that the player is still in the squad after
the round, though the added complexity may not pay for itself. Third,
bench ordering under the auto-substitution rule could be optimized
directly, spreading rotation risk across countries instead of ranking
bench players by predicted points alone.

\subsection{Prediction-Market Combiner}
\label{sec:future-markets}

Section~\ref{sec:negative-results} showed that blending market-implied
signals directly into the expected-points models worsened accuracy.
The follow-up we consider publishable keeps the two inferences
separate and combines them only at the output level with the logit
combiner of \citet{benter1994computer}: the combined forecast is a
logistic regression on the logits of the model forecast and the
market-implied probability, with coefficients fitted on historical
out-of-sample data. Prediction markets aggregate informed and
uninformed money \citep{wolfers2004prediction}; fitting the market
coefficient empirically measures how much signal the market carries
rather than assuming it, and if the market is mostly noise the
coefficient shrinks toward zero.

The implementation plan is concrete. We collected market snapshots
throughout the tournament (152 three-hourly snapshots from the round
of 32 onward, trophy prices for all 50 listed countries); after the final we scrape the closed
contracts for complete price histories, pair them with our archived
per-snapshot predictions and realized outcomes, and fit the combiner
by Bayesian logistic regression, reporting credible intervals that are
necessarily wider than the horse-racing precedent given football's low
scoring and heavy-tailed outcomes. Validation splits within the
tournament (group stage as training, knockout rounds as test) provide
the out-of-sample check. We expect the market to carry most signal on
match outcomes and premium scorers, and our model to retain its edge
on fantasy-specific components such as cheap defenders' clean sheets
and bench mechanics, which no market prices directly. A live variant
also follows: between early and late kickoffs on a matchday, updated
market prices can inform a captain switch among players whose matches
have not yet started.

\subsection{Open Release}
\label{sec:future-release}

The final planned step is the public release of the code, the data
pipeline, and the live decision log paired with realized outcomes,
which accompanies this paper as of 19 July 2026 at
\url{https://github.com/AsherDLL/fifa-wc-fantasy}. Publishing the decision
log lets others reproduce our matchday calls, identify failure modes
we missed, and benchmark alternative predictors against a documented
in-tournament deployment, in line with open-data reproducibility
practice in sports analytics \citep{eager2023football}.
